\begin{lemma}[Weak bound for the square of the local maximal function of a coordinate endpoint function]
Fix $n \geq 1$ and $j \in \{1,\ldots,n\}$, let $\lambda$ be Lebesgue measure on
$\mathbb{R}^n$, let $\gamma_n$ be the Gaussian measure of
\noderef{GaussMeasure}, and let
\[
  T : L^1(\gamma_n;\mathbb{R}) \longrightarrow
      \mathcal{M}\bigl(\gamma_n;\mathbb{R}^n\bigr)
\]
be \emph{any} linear map into the $\gamma_n$-almost-everywhere equivalence
classes of measurable $\mathbb{R}^n$-valued functions such that
\begin{enumerate}
\item[(T1)] $T(\iota v) = R_\gamma v$ for every $v \in L^2(\gamma_n)$, with
$\iota$ the inclusion of \noderef{ToL1} and $R_\gamma$ the transform of
\noderef{RieszTransform}; and
\item[(T2)] $t\,\gamma_n\bigl(\{x : |Tf(x)|_{\mathbb{R}^n} > t\}\bigr)
\leq 2\,\|f\|_{L^1(\gamma_n)}$ for every $f \in L^1(\gamma_n;\mathbb{R})$ and
every real $t > 0$.
\end{enumerate}
Let $R \in \mathbb{R}$, put $\overline B = \overline B(0,R+2)$, let $M_2$ be the
local maximal function of \noderef{LocalMaximalFunction}, and for
$f \in L^1(\gamma_n)$ --- written also for its canonical representative --- put
\[
  u_f = \bigl|(Tf)_j\bigr|^{1/2}\,\mathbf 1_{\overline B}
  \,:\, \mathbb{R}^n \to \mathbb{R} .
\]
Then for every $f \in L^1(\gamma_n)$ and every $s \in (0,+\infty]$,
\[
  \gamma_n\Bigl(\bigl\{\,x \in \mathbb{R}^n :\ M_2(u_f)(x)^2 > s\,\bigr\}\Bigr)
  \;\leq\;
  \frac{32\cdot4^n\,e^{(R+2)^2}\,\|f\|_{L^1(\gamma_n)}}{s} ,
\]
both sides read in $[0,+\infty]$, with $M_2(u_f)(x)^2$ the square taken there.
Here $\|f\|_{L^1(\gamma_n)} = \int_{\mathbb{R}^n}|f|\,d\gamma_n$, the norm of the
class of $f$ written as the integral of the absolute value of its canonical
representative; the two are the same number, and it is that integral the
displayed bound uses.

\emph{Why the square appears in the statement.}  The intended consumer reaches
this term as $\{x : \Psi(x)^2 > s\}$ with $\Psi = M_2(u_f)$, because its own
input bounds a quantity by $\Psi^2 + \Phi + N$ and it splits that sum at the
level $s$; so the squared form is the shape in which the set arises, and
passing to the unsquared level $\{M_2(u_f) > \sqrt s\}$ is part of the work
done here rather than something the consumer must do.  The square is taken in
$[0,+\infty]$: $M_2(u_f)(x)$ is a supremum of averages which is not known in
advance to be finite, exactly as in \noderef{LocalMaximalFunction}.

\emph{The threshold is extended-real, and the endpoint is a triviality.}  At
$s = +\infty$ the level set is empty, nothing in $[0,+\infty]$ exceeding
$+\infty$, while the right-hand side is a finite quantity divided by $+\infty$,
hence $0$; the assertion there is $0 \leq 0$ and the content is at finite $s$.
The threshold is excluded only at $s = 0$, where the right-hand side would be
$+\infty$ unless $\|f\|_{L^1(\gamma_n)} = 0$ and the assertion would fail for
$f$ of norm zero with $M_2(u_f)$ not identically $0$ --- it is $0/0$ that is
being avoided, not a loss of content.

\emph{No measurability of the level set is asserted or needed.}  $\gamma_n$ is
applied to the level set as an outer measure, which is defined on every subset
and is monotone there; and the two facts chained below --- the weak bound of
\noderef{LocalMaximalWeakL2} and the comparison
$\gamma_n \leq \pi^{-n/2}\lambda$ of \noderef{GaussLebesgueBallComparison} ---
are themselves stated for arbitrary sets.  Bounding the outer measure is the
stronger assertion, so a consumer that knows the set to be measurable loses
nothing.

\emph{The constant, and what it does not claim.}  $32\cdot4^n e^{(R+2)^2}$
depends only on $n$ and $R$: the $4^n$ is the Vitali constant carried by
\noderef{LocalMaximalWeakL2}, and $e^{(R+2)^2}$ enters through
$a^{-1} = \pi^{n/2}e^{(R+2)^2}$, the reciprocal of the lower bound
$a = \pi^{-n/2}e^{-(R+2)^2}$ for the Gaussian density on $\overline B$.  No
uniformity in $n$ or $R$ is claimed.  The quantity bounded is coordinatewise in
$j$; nothing here asserts anything about the full vector $Tf$.

\emph{The hypothetical style.}  (T1) and (T2) are written out in full and not
imported by reference; the lemma is a statement about every $T$ satisfying
them, as \noderef{CoordinateEndpointBound} is, and it is not vacuous because
\noderef{L1Extension} exhibits such a $T$.
\end{lemma}
\begin{proof}
Write $p = \pi^{-n/2}$, $q = \pi^{n/2}$ and $a = \pi^{-n/2}e^{-(R+2)^2}$, all in
$(0,\infty)$, so that
\begin{equation*}
  p\,q = 1 \qquad\text{and}\qquad a\cdot\bigl(q\,e^{(R+2)^2}\bigr) = 1 ,
  \tag{1}
\end{equation*}
the first because $\pi^{-n/2} = (\pi^{n/2})^{-1}$ and $\pi^{n/2} \neq 0$, and
the second by multiplying the first by
$e^{-(R+2)^2}e^{(R+2)^2} = e^{-(R+2)^2+(R+2)^2} = e^0 = 1$.  Fix
$f \in L^1(\gamma_n)$, abbreviate $u = u_f$ and $\overline B =
\overline B(0,R+2)$, and put $N = \|f\|_{L^1(\gamma_n)} = \int|f|\,d\gamma_n
\in [0,\infty)$.

\emph{Step 1: $u$ is measurable, nonnegative, and supported in
$\overline B$.}  The canonical representative of a $\gamma_n$-almost-everywhere
equivalence class of measurable functions is measurable, so $(Tf)_j$ is
measurable.  The map $r \mapsto |r|^{1/2}$ is continuous on $\mathbb{R}$, hence
Borel, so $|(Tf)_j|^{1/2}$ is measurable; and $\overline B$ is closed, hence
measurable, so $u = |(Tf)_j|^{1/2}\mathbf 1_{\overline B}$ is measurable.  Since
$|(Tf)_j|^{1/2} \geq 0$ we have $u \geq 0$ everywhere, and $u(y) = 0$ for
$y \notin \overline B$.

\emph{Step 2: the uniform bound on the weak-$L^2$ quantity.}  Put
\[
  \mathcal{A} = \sup_{0<\sigma\leq+\infty}
    \sigma^2\,\lambda\bigl(\{y : |u(y)| > \sigma\}\bigr)
  \;\in\;[0,+\infty] .
\]
We claim
\begin{equation*}
  \mathcal{A} \;\leq\; 2\,q\,e^{(R+2)^2}\,N . \tag{2}
\end{equation*}
It suffices to bound each term of the supremum by the right-hand side, since
that right-hand side does not depend on $\sigma$.  (The uniformity is the
point: a bound still depending on $\sigma$ would not bound the supremum.)

If $\sigma = +\infty$ then $\{y : |u(y)| > \sigma\}$ is empty, $u$ being real
valued so that $|u(y)| \leq +\infty$ always, and the term is
$(+\infty)\cdot 0 = 0$, which is $\leq$ the right-hand side.

So let $\sigma$ be a real number with $\sigma > 0$ and put
$E_\sigma = \{y : |u(y)| > \sigma\}$.  By Step 1, $|u| = u$, so $E_\sigma$ is
measurable, being a strict superlevel set of a measurable function.  Also
$E_\sigma \subseteq \overline B$: if $y \notin \overline B$ then $u(y) = 0$ by
Step 1, so $|u(y)| = 0 \not> \sigma$.  Next,
\begin{equation*}
  E_\sigma \;\subseteq\; \bigl\{y : |(Tf)_j(y)| > \sigma^2\bigr\} . \tag{3}
\end{equation*}
Indeed let $y \in E_\sigma$.  Then $y \in \overline B$ by the previous
sentence, so $u(y) = |(Tf)_j(y)|^{1/2}$, and $|u(y)| = u(y) > \sigma$ gives
$|(Tf)_j(y)|^{1/2} > \sigma > 0$; squaring the two sides of this inequality
between nonnegative reals and using
$\bigl(|(Tf)_j(y)|^{1/2}\bigr)^2 = |(Tf)_j(y)|$, which holds because
$|(Tf)_j(y)| \geq 0$, gives $|(Tf)_j(y)| > \sigma^2$.

Now $E_\sigma$ is measurable and contained in $\overline B = \overline
B(0,R+2)$, so clause (2) of \noderef{GaussLebesgueBallComparison},
\emph{instantiated at $R+2$ in place of its $R$ and at $E := E_\sigma$}, gives
\begin{equation*}
  a\,\lambda(E_\sigma) \;\leq\; \gamma_n(E_\sigma) . \tag{4}
\end{equation*}
And \noderef{CoordinateEndpointBound} --- applied to this $T$ with its two
properties, which are (T1) and (T2) and hold by assumption, to this $f$, to the
coordinate $j$ and to the real level $t := \sigma^2 > 0$ --- gives
\begin{equation*}
  \sigma^2\,\gamma_n\bigl(\bigl\{y : |(Tf)_j(y)| > \sigma^2\bigr\}\bigr)
    \;\leq\; 2\,\|f\|_{L^1(\gamma_n)} \;=\; 2N . \tag{5}
\end{equation*}
Multiplying $(4)$ by $\sigma^2 \in (0,\infty)$, then using monotonicity of
$\gamma_n$ as an outer measure together with $(3)$, then $(5)$:
\[
  a\cdot\sigma^2\lambda(E_\sigma)
    \;\leq\; \sigma^2\gamma_n(E_\sigma)
    \;\leq\; \sigma^2\gamma_n\bigl(\bigl\{y : |(Tf)_j(y)| > \sigma^2\bigr\}\bigr)
    \;\leq\; 2N
    \;=\; a\cdot\bigl(2\,q\,e^{(R+2)^2}N\bigr) ,
\]
the last equality being $(1)$ multiplied by $2N$.  Cancelling the factor
$a \in (0,\infty)$, which is legitimate in $[0,+\infty]$ because $a$ is neither
$0$ nor $+\infty$, gives
$\sigma^2\lambda(E_\sigma) \leq 2\,q\,e^{(R+2)^2}N$.  As this holds for every
$\sigma$, $(2)$ follows.

\emph{Step 3: the level $s$, its square root, and the conclusion.}  Let
$s \in (0,+\infty]$ and put
$S = \{x : M_2(u)(x)^2 > s\}$.

If $s = +\infty$ then $S$ is empty, since $M_2(u)(x)^2 \leq +\infty$ for every
$x$, so $\gamma_n(S) = 0$ and the assertion holds, the right-hand side being
$\geq 0$.

So let $s \in (0,\infty)$ and let $\sqrt s \in (0,\infty)$ be its positive
square root, so $(\sqrt s)^2 = s$ and $\sqrt s \neq 0$.  Then
\begin{equation*}
  S \;\subseteq\; \bigl\{x : M_2(u)(x) > \sqrt s\bigr\} , \tag{6}
\end{equation*}
by contraposition: if $M_2(u)(x) \leq \sqrt s$ then, squaring is monotone on
$[0,+\infty]$, $M_2(u)(x)^2 \leq (\sqrt s)^2 = s$, so $x \notin S$.

Apply \noderef{LocalMaximalWeakL2} to $u$, which is measurable hence measurable
up to a null set by Step 1, at the level $\sqrt s \neq 0$.  Together with
$(\sqrt s)^2 = s$ it gives
\[
  \lambda\bigl(\bigl\{x : M_2(u)(x) > \sqrt s\bigr\}\bigr)
    \;\leq\; \frac{16\cdot4^n\,\mathcal{A}}{s} .
\]
By monotonicity of $\lambda$ as an outer measure and $(6)$, the same bound
holds for $\lambda(S)$.  By clause (1) of
\noderef{GaussLebesgueBallComparison}, which holds for \emph{every} set,
$\gamma_n(S) \leq p\,\lambda(S)$.  Combining these with $(2)$ and using
monotonicity of $r \mapsto p\cdot 16\cdot 4^n r/s$ on $[0,+\infty]$,
\[
  \gamma_n(S)
    \;\leq\; p\,\lambda(S)
    \;\leq\; \frac{p\cdot16\cdot4^n\,\mathcal{A}}{s}
    \;\leq\; \frac{p\cdot16\cdot4^n\cdot2\,q\,e^{(R+2)^2}N}{s}
    \;=\; \frac{32\cdot4^n\,e^{(R+2)^2}N}{s} ,
\]
the last step because
$p\cdot16\cdot4^n\cdot2\,q\,e^{(R+2)^2}N
 = 32\cdot4^n\,e^{(R+2)^2}N\cdot(p\,q)$ and $p\,q = 1$ by $(1)$.  This is the
assertion for $s$, and with the case $s = +\infty$ it holds for every
$s \in (0,+\infty]$.
\end{proof}
